%% sections/02-related-work.tex

\section{Related Work}
\label{sec:related}

\subsection{Rubric Evolution in Educational Assessment}

The educational assessment literature has long recognized that rubrics require
revision when their anchor descriptions fail to capture the range of student work
encountered in practice~\cite{andrade2000, nitko2001}. Revision processes range from
informal annual updates to structured validity studies comparing rubric scores with
expert judgments~\cite{boud2013assessment}. The Marathon mechanism formalizes a
bottom-up version of this process: instead of expert judgment as the revision trigger,
it uses systematic logging of evaluator surprises at score time.

\subsection{Active Learning and Annotation}

Active learning addresses the analogous problem in machine learning annotation: new
data points may fall outside the current label set~\cite{settles2009active}. When
annotators encounter out-of-vocabulary examples, they either force-map to the nearest
existing category (introducing noise) or flag the example for label expansion. The
Marathon mechanism is a structured version of the flag-and-expand path: innovations
are flagged systematically; expansion happens when flagging density exceeds a threshold.

This is related to but distinct from \textit{ontology evolution}~\cite{noy2001ontology},
in which a knowledge representation is extended as new concepts are encountered in
domain practice. Rubric amendment is domain-constrained (the rubric scores a specific
type of artifact --- AI-simulated TRPG sessions --- rather than a general domain) and
operationally motivated (amendments are driven by evaluator surprises, not domain
coverage requirements).

\subsection{Iterative Refinement in Evaluation Systems}

Several evaluation systems for generative AI use iterative refinement mechanisms.
\citet{zheng2023judging} describe MT-Bench as a living benchmark updated when
new model capabilities escape existing question categories. \citet{liang2022holistic}
describe HELM's dimension set as subject to expansion as the field evolves.
Both systems use expert review to drive expansion rather than systematic logging of
evaluator surprises. The Marathon mechanism is more automated: the innovation log
is part of the evaluation pipeline, not a separate review cycle.

\subsection{Scoring System Drift}

A concern with iterative rubric revision is \textit{scoring drift}: scores from
different rubric versions may not be comparable~\cite{angoff1971scales}. The
forward-only constraint in the Marathon mechanism directly addresses this: scores
from v1.3 sessions are comparable to each other, and cross-version comparisons
are explicit rather than conflated. This is analogous to semantic versioning in
software development~\cite{semver}, where backward-incompatible changes increment
the major version and users know not to compare across major versions.
